\section{The massless limit}
\label{sec:massless}

In this section we take the limit $\mf\to0$ of our results and compare with the known expressions for
massless quarks~\cite{Iancu:2022gpw}.

\subsection{Scales, vertices and radial functions}

For $\mf\to0$ the mass scales of eqs.~\eqref{eq:notation} and \eqref{eq:notation-q} become
\begin{equation}
\ef^2\to\bar Q^2\equiv\zeta(1-\zeta)\,Q^2\,,\quad\;\eft^{\,2}\to\widetilde Q^2\equiv\eta(1-\eta)\,Q^2\,,\quad\;
\of,\,\oft\to0\,,\quad\;\Ef\to Q\,.
\end{equation}
The helicity-flip terms of the two vertices \eqref{eq:vertices} vanish, so that the helicity is
conserved along the fermion line and only the first of the four helicity paths of
eq.~\eqref{eq:Psi} survives. This requires the products $\mf\,\cG^{(01)}$, $\xi\of\,\cG^{(10)}$ and
$\mf\,\xi\of\,\cG^{(00)}$ to vanish, which is the case because the three functions have finite limits
for $Q^2>0$, and similarly for the functions $\cA^{(\alpha\beta)}$, for which the worst behaviour is
$\xi\of K_0(\of Y)\sim\of\ln(1/\of)\to0$.

The surviving function $\cG^{(11)}$ can be computed in closed form. For $\mf=0$ one has
$A_u=\bar Q^2$ independently of $u$, and with the change of variable $s=\Delta_u=\R^2+a\Y^2/u$,
\begin{equation}
\cG^{(11)}\Big|_{\mf=0}=\frac{1}{2\zeta(1-\zeta)\,\Y^2}\int_{\R^2+a\Y^2}^\infty\frac{\dd s}{s}\;\bar Q^2\,
K_2\big(\bar Q\sqrt s\big)
=\frac{1}{\zeta(1-\zeta)\,\Y^2}\;\frac{\bar Q\,K_1\big(\bar Q\sqrt{\R^2+a\Y^2}\big)}{\sqrt{\R^2+a\Y^2}}\,,
\end{equation}
where we used $\int_t^\infty\dd t'\,K_2(t')/t'=K_1(t)/t$. Since $\zeta(1-\zeta)(\R^2+a\Y^2)=\cD^2$, cf.\
eq.~\eqref{eq:D}, this is
\begin{equation}
\cG^{(11)}\to\frac{Q\,K_1(Q\cD)}{\cD\;\Y^2}\,,\qquad
\cA^{(11)}\to\frac{\bar Q\,K_1(\bar QR)}{\zeta(1-\zeta)\,R\;\Y^2}=\frac{Q^2K_1(\bar QR)}{\bar QR\;\Y^2}\,,\qquad
\cG_E\to\frac{Q\,K_1(Q\cD)}{\cD}\,,
\label{eq:massless-fns}
\end{equation}
in agreement with eq.~\eqref{eq:FT-massless}. In this limit $\cA^{(11)}$ is the value of $\cG^{(11)}$ at
$\y,\z\to\yp$ at fixed $\R$ and with the factor $1/\Y^2$ kept apart, since $\cD\to\sqrt{\zeta(1-\zeta)}\,R$
there: this is the origin of the formal substitution rule of ref.~\cite{Iancu:2022gpw}, which has no
simple analogue in the massive case. The effective vertices \eqref{eq:M} become
\begin{align}
\cM^{im}_{1\,\lambda_1\lambda_2}&\to\bigg[\varphi^{ij}_{\lambda_1\lambda}(\zeta)\,\tau^{mn}_{\lambda\lambda_2}(\xi,\eta)\,
\frac{R^jY^n}{\Y^2}-\delta_{\lambda_1\lambda_2}\,\frac{\xi\,\eta}{1-\zeta}\,\big(\delta^{im}+2i\ee^{im}\lambda_1\big)\bigg]
\frac{Q\,K_1(Q\cD)}{\cD}\,,
\nonumber\\
\cM^{im}_{2\,\lambda_1\lambda_2}&\to\varphi^{ij}_{\lambda_1\lambda}(\zeta)\,\tau^{mn}_{\lambda\lambda_2}(\xi,\eta)\,
\frac{R^jY^n}{\Y^2}\;\frac{Q^2K_1(\bar QR)}{\bar QR}\,,
\label{eq:M-massless}
\end{align}
and the outgoing state \eqref{eq:state-qbar} reduces to eqs.~(4.25), (4.28) and (4.29) of
ref.~\cite{Iancu:2022gpw}. For the emission by the quark,
\begin{equation}
\tG^{(11)}\to\frac{Q\,K_1(Q\cD)}{\cD\;\X^2}\,,\qquad
\tA^{(11)}\to\frac{Q^2K_1(\widetilde Q\tilde R)}{\widetilde Q\tilde R\;\X^2}\,,
\end{equation}
and eq.~\eqref{eq:state-q} reduces to eqs.~(4.33)--(4.35) of the same reference.

\subsection{Cross-section}

In the helicity sums \eqref{eq:H-direct}, \eqref{eq:Ht-direct} and \eqref{eq:Hx} all the terms
proportional to $\mf$ drop out and we are left with
\begin{align}
\cH\big|_{\mf=0}={}&\Big\{8\big[\zeta^2+(1-\zeta)^2\big]\big[\eta^2+(1-\zeta)^2\big]\,(\R\cdot\bR)(\Y\cdot\bY)
\nonumber\\
&\quad-8\,\xi\,(2\zeta-1)\big(2(1-\zeta)-\xi\big)\,(\R\times\bR)(\Y\times\bY)\Big\}\,F^{(11)}\bar F^{(11)}
\nonumber\\
&-\frac{8\,\zeta\,\xi\,\eta^2}{1-\zeta}\,\Big[(\R\cdot\Y)\,F^{(11)}\bar F_E+(\bR\cdot\bY)\,F_E\,\bar F^{(11)}\Big]
+8\left[\frac{\xi\,\eta}{1-\zeta}\right]^2F_E\,\bar F_E\,,
\nonumber\\[1mm]
\tH\big|_{\mf=0}={}&\Big\{8\big[\zeta^2+(1-\eta)^2\big]\big[\eta^2+(1-\eta)^2\big]\,(\Rt\cdot\bRt)(\X\cdot\bX)
\nonumber\\
&\quad-8\,\xi\,(2\eta-1)\big(2(1-\eta)-\xi\big)\,(\Rt\times\bRt)(\X\times\bX)\Big\}\,\tF^{(11)}\bar\tF^{(11)}
\nonumber\\
&+\frac{8\,\zeta^2\,\xi\,\eta}{1-\eta}\,\Big[(\Rt\cdot\X)\,\tF^{(11)}\bar\tF_E+(\bRt\cdot\bX)\,\tF_E\,\bar\tF^{(11)}\Big]
+8\left[\frac{\xi\,\zeta}{1-\eta}\right]^2\tF_E\,\bar\tF_E\,,
\nonumber\\[1mm]
\cH^\times\big|_{\mf=0}={}&\Big\{8\big[1-\xi-2\zeta\eta\big]\big[\xi(1-\xi)+2\zeta\eta\big]\,(\Rt\cdot\bR)(\X\cdot\bY)
\nonumber\\
&\quad-8\,\xi\,\big[1-2\zeta-\xi\big]^2\,(\Rt\times\bR)(\X\times\bY)\Big\}\,\tF^{(11)}\bar F^{(11)}
\nonumber\\
&-\frac{8\,\xi\,\eta\,(\zeta+\xi)^2}{1-\zeta}\,(\Rt\cdot\X)\,\tF^{(11)}\bar F_E
+\frac{8\,\xi\,\zeta\,(1-\zeta)^2}{\zeta+\xi}\,(\bR\cdot\bY)\,\tF_E\,\bar F^{(11)}\,,
\label{eq:H-massless}
\end{align}
with the functions \eqref{eq:massless-fns}. The regular~$\times$~regular parts reproduce the squared
vertices (5.5) and (5.17) of ref.~\cite{Iancu:2022gpw}, and the terms involving $F_E$ are the explicit
form of the instantaneous contributions, which in that reference are kept inside the effective
vertices. For instance, the term of eq.~\eqref{eq:xs-qbar} in which the three partons scatter on both
sides of the cut becomes
\begin{align}
\frac{\dd\sigma^{(\bar q\bar q)}}{\dd^3k_1\,\dd^3k_2\,\dd^3k_3}\bigg|_{\mf=0}\ni{}&
\frac{\aem\,\as\,\CF\Nc}{2(2\pi)^{10}}\;\frac{\zeta^2\,Q^2}{\xi\,(q^+)^2}\;
\delta\big(q^+-k_1^+-k_2^+-k_3^+\big)\sum_fe_f^2
\nonumber\\
&\times\int_{\x,\y,\z,\bx,\by,\bz}e^{-i\kk_1\cdot(\x-\bx)-i\kk_2\cdot(\y-\by)-i\kk_3\cdot(\z-\bz)}\;
\frac{K_1(Q\cD)\,K_1(Q\bar\cD)}{\cD\,\bar\cD}
\nonumber\\
&\times\bigg\{\frac{8\big[\zeta^2+(1-\zeta)^2\big]\big[\eta^2+(1-\zeta)^2\big](\R\cdot\bR)(\Y\cdot\bY)}{\Y^2\,\bY^2}
\nonumber\\
&\qquad-\frac{8\,\xi\,(2\zeta-1)\big(2(1-\zeta)-\xi\big)(\R\times\bR)(\Y\times\bY)}{\Y^2\,\bY^2}
\nonumber\\
&\qquad-\frac{8\,\zeta\,\xi\,\eta^2}{1-\zeta}\left[\frac{\R\cdot\Y}{\Y^2}+\frac{\bR\cdot\bY}{\bY^2}\right]
+8\left[\frac{\xi\,\eta}{1-\zeta}\right]^2\bigg\}\;\cW(\x,\y,\z;\bx,\by,\bz)\,,
\end{align}
which is eq.~(5.2) of ref.~\cite{Iancu:2022gpw}. In momentum space, the massless limit of
section~\ref{sec:mom} is obtained by dropping $H^{\rm massive}$, $\widetilde H^{\rm massive}$ and
$H_\times^{\rm massive}$ and by setting $\ef\to\bar Q$ and $\of\to0$ in eq.~\eqref{eq:f12},
\begin{align}
f_1&\to\frac{1-\zeta}{\big(\kk^2+\bar Q^2\big)\Big[\zeta(1-\zeta)^2\pp^2+\xi\eta\big(\kk^2+\bar Q^2\big)\Big]}\,,
\nonumber\\
f_2&\to\frac{1}{\zeta(1-\zeta)\big(\kk^2+\bar Q^2\big)\,\pp^2}\,,\qquad
g\to\frac{1}{\zeta(1-\zeta)^2\pp^2+\xi\eta\big(\kk^2+\bar Q^2\big)}\,.
\label{eq:f-massless}
\end{align}

\subsection{How the limit is approached}

Two remarks are in order about the nature of the limit. First, at fixed transverse coordinates the
corrections to the massless results are of relative order $\mf^2/\bar Q^2$ for the helicity-conserving
structures, and the helicity-flip terms contribute to the cross-section at order $\mf^2$. We have
checked numerically that $\cG^{(11)}$ approaches eq.~\eqref{eq:massless-fns} quadratically in $\mf$ and
that the flip terms vanish linearly (appendix~\ref{app:checks}).

Second, the limit is not uniform over the phase space of the gluon. For an emission after the
collision the massive kernel $\of K_1(\of Y)/Y$ reduces to the massless one, $1/\Y^2$, only for
$Y\ll1/\of$, whereas it is exponentially suppressed for $Y\gtrsim1/\of$. Equivalently, in momentum
space $f_2\propto1/(\pp^2+\of^2)$ is cut off at relative momenta $|\pp|\lesssim\of$. The massless limit
therefore does not commute with the integration over the gluon phase space: the cross-section
integrated over the gluon contains a logarithm $\ln(\mu^2/\mf^2)$, which turns into the collinear
divergence of the massless calculation. We come back to this point in section~\ref{sec:coll}.
Similarly, the massless limit does not commute with the photoproduction limit: for $Q^2\to0$ the size
of the pair is controlled by $\ef\to\mf$ and the limit $\mf\to0$ would produce an infrared divergence
at large dipole sizes, which in reality is regulated by confinement.
